\documentclass{article}
\usepackage[utf8]{inputenc}
\usepackage[english]{babel}
\usepackage{geometry}
\usepackage{amsmath}
\usepackage{graphicx}

\geometry{margin=1in, a4paper}

\title{PSVO Zadanie 1 - dokuementácua}
\author{Tadeáš Kapina, Radoslav Janíček}
\date{\today}

\begin{document}

\maketitle

\section{Kalibrácia Kamery}

Kalibrácia kamery prebehla podľa oficiálneho návodu OpenCV. Nasnímali sme 22 snímok šachovnice, na ktorých prebehla kalibrácia. 
Snímky sme nasnímali pred samotným kalibračným procesom a uložili ich do priečinka - snímky sme ukladali iba keď na nich OpenCV dokázal rozoznať šachovnicu. 

Následne sme použili funkciu \texttt{cv.findChessboardCorners()} na nájdenie rohov a následne \texttt{cv.cornerSubPix()}, ktorá slúži na spresnenie detektie.

Nakoniec sme použili funkciu \texttt{cv2.calibrateCamera()} na získanie matice vnútorných parametrov kamery a disorčných koeficientov.

\subsection{Parametre kamery}

\subsubsection*{Matica vnútorných parametrov kamery}
\[
K = \begin{pmatrix}
600.352 & 0 & 182.412 \\
0 & 722.178 & 192.004 \\
0 & 0 & 1
\end{pmatrix}
\]

\subsubsection*{Vnútorné parametre kamery}
\begin{itemize}
    \item $f_x = 600.352$
    \item $f_y = 722.178$
    \item $c_x = 182.412$
    \item $c_y = 192.004$
\end{itemize}

\subsubsection*{Disorčné koeficienty}
\[
D = \begin{pmatrix}
-0.399 & -0.417 & 0.001 & 0.003 & 3.403
\end{pmatrix}
\]

\section{Detekcia geometrických tvarov}


\subsection{Použité metódy}

Na detekciu základných geometrických tvarov v obraze boli využité metódy a funkcie z knižnice OpenCV:

\begin{itemize}
    \item \textbf{Detekcia kružníc:} Využíva sa Houghova transformácia na hľadanie kruhov (\texttt{cv2.HoughCircles}). Metóda sa aplikuje na obraz, ktorý bol najprv prevedený do odtieňov sivej a vyhladený mediánovým filtrom pre zníženie šumu.
    \item \textbf{Detekcia polygónov (štvorce, obdĺžniky, trojuholníky):}
    \begin{itemize}
        \item Na nájdenie hrán v obraze je použitý Cannyho detektor hrán (\texttt{cv2.Canny}).
        \item Následne program hľadá spojené oblasti (kontúry) pomocou funkcie \texttt{cv2.findContours}.
        \item Nájdené kontúry sú aproximované pomocou \texttt{cv2.approxPolyDP} na zistenie počtu vrcholov (napríklad 3 vrcholy pre trojuholník, 4 pre štvorec alebo obdĺžnik).
    \end{itemize}
    \item \textbf{Určenie ťažiska:} Stred (ťažisko) každého detegovaného tvaru sa vypočítava z obrazových momentov pomocou funkcie \texttt{cv2.moments}.
\end{itemize}

\subsection{Postup riešenia}

Postup spracovania obrazu a detekcie tvarov funguje nasledovne:

\begin{enumerate}
    \item Program v nekonečnej slučke načítava obraz z kamery.
    \item Snímka sa konvertuje do grayscale formátu a filtruje sa, aby sa pripravila na detekciu.
    \item Pre zabezpečenie robustnosti voči rôznym svetelným podmienkam a scénam program obsahuje grafické rozhranie s posuvníkmi (Trackbars). Tieto umožňujú používateľovi v reálnom čase ladiť prahy pre Cannyho detektor, parametre Houghovej transformácie, úroveň aproximácie a minimálnu plochu objektov.
    \item Ak program deteguje tvar, graficky ho v reálnom obraze označí (vykreslí obrys zelenou farbou).
    \item Do stredu tvaru sa vykreslí červený bod, ktorý reprezentuje jeho ťažisko.
    \item Ku každému rozpoznanému tvaru program pridá textovú menovku s jeho názvom v reálnom čase (napr. \texttt{Tri}, \texttt{Square}, \texttt{Rect}, \texttt{Circle}).
\end{enumerate}


\section{Farebný filter}

Pre zmenu jednej farby na iný sme najskr zmenili farebný priestor z RGB na HSV pomocou funkcie \texttt{cv2.cvtColor}. Následne sme použili funkciu \texttt{cv2.inRange} na vytvorenie masky, ktorá vyberá pixely v zadanom rozsahu farieb v HSV priestore. Nakoniec sme všekty pixely, ktoré patrili do masky zmenili na požadovanú farbu.

\end{document}